\documentclass[12pt,a4paper]{article}

\usepackage[T1]{fontenc}
\usepackage[utf8]{inputenc}
\usepackage[french]{babel}
\usepackage{geometry}
\usepackage{lmodern}
\usepackage{setspace}
\usepackage{hyperref}
\usepackage{longtable}
\usepackage{array}
\usepackage{enumitem}
\usepackage{xcolor}

\geometry{margin=2.5cm}
\onehalfspacing
\hypersetup{colorlinks=true,linkcolor=blue,urlcolor=blue}

\title{Rapport Technique\\\large Systeme de Gestion des Auditoires et des Horaires}
\author{Projet PHP Procedural (sans framework)}
\date{\today}

\begin{document}
\maketitle
\tableofcontents
\newpage

\section{Introduction}
Ce document presente de maniere detaillee le travail realise pour concevoir et developper une application complete de gestion des auditoires.
L'application permet de planifier les cours, reserver les salles, eviter les conflits d'occupation et optimiser l'affectation des ressources selon la capacite des auditoires.

Le projet est construit en \textbf{PHP procedural} (sans framework), avec des vues HTML/CSS integrees dans les fichiers PHP et un stockage des donnees en \textbf{JSON}.

\section{Objectifs fonctionnels}
\subsection{Objectif general}
Mettre en place un systeme academique permettant de:
\begin{itemize}[leftmargin=1.2cm]
    \item planifier les cours;
    \item reserver les salles;
    \item attribuer les salles selon la capacite des groupes;
    \item bloquer les conflits (meme salle, meme creneau);
    \item rendre le planning coherent, lisible et exploitable.
\end{itemize}

\subsection{Contraintes respectees}
\begin{itemize}[leftmargin=1.2cm]
    \item aucun framework PHP;
    \item persistance uniquement via des fichiers JSON;
    \item logique metier centralisee dans des fonctions reutilisables;
    \item interfaces admin et publiques distinctes;
    \item controle d'acces administrateur avec gestion de droits.
\end{itemize}

\section{Architecture du projet}
\subsection{Arborescence logique actuelle}
\begin{itemize}[leftmargin=1.2cm]
    \item \texttt{index.php}: page publique principale (accueil + apercus);
    \item \texttt{config/config.php}: configuration globale de l'application;
    \item \texttt{data/*.json}: base de donnees JSON;
    \item \texttt{includes/}: fonctions, header, footer;
    \item \texttt{admin/}: espace administrateur (CRUD + planning + securite + gestion admins);
    \item \texttt{actions/}: endpoints de traitement (ajout, generation, reset, etc.);
    \item \texttt{auth/login.php} et \texttt{auth/logout.php}: authentification;
    \item \texttt{pages/stats.php}: statistiques et disponibilites;
    \item \texttt{assets/css/style.php}, \texttt{assets/js/app.js}: ressources front.
\end{itemize}

\subsection{Separation des responsabilites}
\begin{itemize}[leftmargin=1.2cm]
    \item \textbf{Config}: constantes application + resolution d'URL;
    \item \textbf{Fonctions metier}: lecture/criture JSON, validations, planning, permissions;
    \item \textbf{Pages}: presentation et formulaires;
    \item \textbf{Actions}: traitement des soumissions et redirections.
\end{itemize}

\section{Modele de donnees JSON}
\subsection{Fichiers utilises}
\begin{longtable}{|>{\raggedright\arraybackslash}p{3.2cm}|>{\raggedright\arraybackslash}p{9.2cm}|}
\hline
\textbf{Fichier} & \textbf{Contenu} \\\hline
\texttt{salles.json} & id, designation, capacite \\\hline
\texttt{promotions.json} & id\_promotion, libelle, effectif\_total \\\hline
\texttt{options.json} & id\_option, libelle, promotion\_parente, effectif \\\hline
\texttt{cours.json} & id\_cours, intitule, volume\_horaire, promotion\_option \\\hline
\texttt{planning.json} & creneau, salle, cours, groupe \\\hline
\texttt{admin.json} & liste des administrateurs, hash de mot de passe, statut actif, permissions \\\hline
\end{longtable}

\subsection{Relations fonctionnelles}
\begin{itemize}[leftmargin=1.2cm]
    \item un \textbf{cours} cible soit une promotion, soit une option (champ \texttt{promotion\_option});
    \item une \textbf{option} depend d'une promotion parente;
    \item une ligne de \textbf{planning} relie un creneau, une salle, un cours et un groupe.
\end{itemize}

\section{Logique metier principale}
\subsection{Regles de reservation}
Lors de l'ajout (manuel ou automatique) d'une seance:
\begin{itemize}[leftmargin=1.2cm]
    \item tous les champs doivent etre remplis;
    \item le groupe doit exister;
    \item la salle doit etre libre sur le creneau;
    \item le groupe ne doit pas deja etre occupe sur le meme creneau;
    \item la capacite de la salle doit etre suffisante.
\end{itemize}

\subsection{Gestion des conflits}
Deux types de conflits sont verifies:
\begin{itemize}[leftmargin=1.2cm]
    \item \textbf{conflit salle}: meme salle + meme creneau;
    \item \textbf{conflit groupe}: meme groupe + meme creneau.
\end{itemize}

\subsection{Attribution intelligente des salles}
L'algorithme trie les salles par capacite croissante et prend la premiere salle disponible capable d'accueillir le groupe.
Cela minimise le gaspillage des grandes salles et reserve les capacites elevees aux grands effectifs.

\subsection{Generation automatique du planning}
La generation suit une logique en plusieurs etapes:
\begin{enumerate}[leftmargin=1.2cm]
    \item tri des cours par taille de groupe (du plus grand au plus petit), puis volume horaire;
    \item decoupage du volume horaire en sessions (par bloc d'environ 2h);
    \item recherche d'un creneau sans conflit;
    \item recherche de la meilleure salle disponible;
    \item ecriture des seances dans \texttt{planning.json};
    \item retour des erreurs si certaines sessions sont impossibles a placer.
\end{enumerate}

\section{Interfaces et parcours utilisateur}
\subsection{Espace public}
\begin{itemize}[leftmargin=1.2cm]
    \item page d'accueil: KPIs, planning du jour, planning global;
    \item page statistiques: disponibilite des salles, occupation par jour, capacites globales.
\end{itemize}

\subsection{Espace administrateur}
\begin{itemize}[leftmargin=1.2cm]
    \item dashboard (indicateurs de synthese);
    \item CRUD salles, cours, promotions, options;
    \item planning: filtres, vue hebdomadaire, ajout, modification, suppression, generation auto, reset;
    \item changement de mot de passe;
    \item gestion des administrateurs et des droits.
\end{itemize}

\section{Securite et controle d'acces}
\subsection{Authentification}
\begin{itemize}[leftmargin=1.2cm]
    \item connexion via \texttt{auth/login.php};
    \item deconnexion via \texttt{auth/logout.php};
    \item mots de passe stockes en hash (\texttt{password\_hash}/\texttt{password\_verify});
    \item session PHP pour maintenir l'etat authentifie.
\end{itemize}

\subsection{RBAC (Role-Based Access Control)}
Chaque administrateur possede un ensemble de permissions booleennes:
\begin{itemize}[leftmargin=1.2cm]
    \item \texttt{manage\_rooms}
    \item \texttt{manage\_courses}
    \item \texttt{manage\_promotions}
    \item \texttt{manage\_options}
    \item \texttt{manage\_planning}
    \item \texttt{manage\_admins}
\end{itemize}

Ces permissions sont appliquees a la fois sur:
\begin{itemize}[leftmargin=1.2cm]
    \item l'affichage des menus;
    \item l'acces aux pages;
    \item l'execution des actions serveur.
\end{itemize}

\subsection{Regle de securite specifique}
Un administrateur connecte ne peut pas se retirer lui-meme le droit \texttt{manage\_admins}. Cette protection evite de perdre la capacite de gouvernance des comptes administrateurs.

\section{Qualite et robustesse}
\subsection{Validation des donnees}
La validation est appliquee sur les champs critiques (vide, types, cohérence metier).
Les erreurs sont restituees a l'utilisateur via des messages flash clairs.

\subsection{Integrite des ecritures}
Les suppressions/modifications de planning sont basees sur des cles logiques (creneau+salle+cours+groupe), et non sur des index visuels, ce qui reduit les risques d'erreur avec les filtres.

\subsection{Verification technique}
Les fichiers PHP ont ete verifies regulierement avec \texttt{php -l} pour detecter rapidement les erreurs de syntaxe durant l'evolution du code.

\section{Limites actuelles}
\begin{itemize}[leftmargin=1.2cm]
    \item stockage JSON adapte aux petites/moyennes charges mais moins adapte au multi-utilisateur intensif;
    \item absence de transactions fortes comme dans une base relationnelle;
    \item historique/audit des actions administrateur non encore implemente.
\end{itemize}

\section{Perspectives d'amelioration}
\begin{itemize}[leftmargin=1.2cm]
    \item journalisation des actions sensibles (audit trail);
    \item export PDF/Excel du planning;
    \item notification des conflits restants;
    \item moteur de generation plus avance (contraintes pedagogiques supplementaires);
    \item migration optionnelle vers une base SQL si le volume augmente.
\end{itemize}

\section{Conclusion}
Le projet repond aux exigences centrales: planification, reservation, prevention des conflits, respect des capacites, optimisation de l'affectation et administration securisee.
La conception en PHP procedural et JSON est respectee, avec une architecture claire, maintenable et evolutive.

Ce travail constitue une base operationnelle solide pour la gestion academique des auditoires et des horaires, tout en laissant de la place a des extensions futures.

\end{document}
